\subsection{Simplifying Fractions}
There are \makeblank to write any rational number as a fraction.
\begin{definition}[Equivalent Fractions]
Two fractions $\frac{a}{b}$ and $\frac{c}{d}$ are \term{equivalent} if
$$
ad=bc
$$
The products $ad$ and $bc$ are called the \term{cross-products}. If $\frac{a}{b}$ and $\frac{c}{d}$ are equivalent, we write
$$
\tfrac{a}{b}=\tfrac{c}{d}
$$
and think of the fractions as representing the same rational number.
\end{definition}
This definition helps us check whether fractions are equivalent, but doesn't help us get equivalent fractions.
\begin{itemize}
\item For any \makeblank[1in] $c\neq\makeblanksmall$, $\frac{a}{b}=\frac{a\cdot c}{b\cdot c}$
\item For any \makeblank[1in] $c$ such that $a\div c$ and $b\div c$ are both \makeblank[1in] (that is, $c$ is a \makeblank of $a$ and $b$), $\frac{a}{b}=\frac{a\div c}{b\div c}$.
\end{itemize}
\begin{definition}[Lowest Terms]
A fraction $\frac{n}{d}$ is said to be in \term{lowest terms} if
\begin{itemize}
\item $n$ and $d$ have no common factors except \makeblanksmall and \makeblanksmall
\item \makeblanksmall is positive
\end{itemize}
Every fraction is equivalent to \makeblank fraction in lowest terms. The process of putting a fraction in lowest terms is called \term{simplifying} or \term{reducing} the fraction.
\end{definition}
\begin{example}
Simplify the fraction $\frac{-18}{-42}$
\begin{cpic}{}{5}
\end{cpic}
\end{example}
If we use a \makeblank for fraction arithmetic, it automatically gives the result in lowest terms.